% from ega4-2-fr.tex
\oldpage[IV-2]{8}
\begin{corollary}[2.1.10]
\label{IV.2.1.10-fr}
Soient $f:X\to Y$ un morphisme de préschémas, $\mathscr F$ un
$\mathcal O_X$-Module quasi-cohérent et $Y$-plat,
$\mathscr L_\bullet=(\mathscr L_i)$ une résolution gauche de $\mathscr F$
formée de $\mathcal O_X$-Modules quasi-cohérents et $Y$-plats. Alors, pour
tout morphisme $g:Y'\to Y$ et tout $\mathcal O_{Y'}$-Module quasi-cohérent
$\mathscr G'$, le complexe
$\mathscr L_\bullet\otimes_Y\mathscr G'=
(\mathscr L_i\otimes_Y\mathscr G')_{i\geq0}$ est une résolution gauche de
$\mathscr F\otimes_Y\mathscr G'$.

En outre, si
$\mathscr Z_i(\mathscr L_\bullet)=\operatorname{Ker}
(\mathscr L_i\to\mathscr L_{i-1})$, les $\mathscr Z_i(\mathscr L_\bullet)$
sont $Y$-plats, et l'on a
$\mathscr Z_i(\mathscr L_\bullet)\otimes_Y\mathscr G'
=\mathscr Z_i(\mathscr L_\bullet\otimes_Y\mathscr G')
=\operatorname{Ker}(\mathscr L_i\otimes_Y\mathscr G'
\to\mathscr L_{i-1}\otimes_Y\mathscr G')$.
\end{corollary}

Posons
$\mathscr R_i=\operatorname{Im}(\mathscr L_{i+1}\to\mathscr L_i)
=\mathscr Z_i(\mathscr L_\bullet)$~; on a alors les suites exactes
\[
0\leftarrow\mathscr F\leftarrow\mathscr L_0\leftarrow\mathscr R_0\leftarrow0
\]
\[
\hbox to 12em{\dotfill}
\]
\[
0\leftarrow\mathscr R_i\leftarrow\mathscr L_{i+1}
\leftarrow\mathscr R_{i+1}\leftarrow0
\]
\[
\hbox to 12em{\dotfill}
\]
et comme $\mathscr F$ et les $\mathscr L_i$ sont $Y$-plats, on déduit de
(2.1.8, (ii)) par récurrence que tous les $\mathscr R_i$ sont aussi
$Y$-plats~; utilisant (2.1.8, (i)), on a donc les suites exactes
\[
0\leftarrow\mathscr F\otimes_Y\mathscr G'
\leftarrow\mathscr L_0\otimes_Y\mathscr G'
\leftarrow\mathscr R_0\otimes_Y\mathscr G'\leftarrow0
\]
\[
0\leftarrow\mathscr R_i\otimes_Y\mathscr G'
\leftarrow\mathscr L_{i+1}\otimes_Y\mathscr G'
\leftarrow\mathscr R_{i+1}\otimes_Y\mathscr G'\leftarrow0
\qquad(i\geq0)
\]
qui prouvent le corollaire.

\begin{proposition}[2.1.11]
\label{IV.2.1.11-fr}
Soient $f:X\to Y$ un morphisme plat, $\mathscr F$ un $\mathcal O_Y$-Module
quasi-cohérent de présentation finie. Si $\mathscr J$ est l'Idéal de
$\mathcal O_Y$ annulateur de $\mathscr F$, $f^*(\mathscr J)$ est l'Idéal de
$\mathcal O_X$ annulateur de $f^*(\mathscr F)$.
\end{proposition}

On a en effet par définition une suite exacte $(\mathbf{0}_{\mathrm I},5.3.7)$
\[
0\to\mathscr J\to\mathcal O_Y\to
\mathscr H\!om_{\mathcal O_Y}(\mathscr F,\mathscr F)
\]
d'où, puisque $f$ est plat, une suite exacte
\[
0\to f^*(\mathscr J)\to\mathcal O_X\to
f^*(\mathscr H\!om_{\mathcal O_Y}(\mathscr F,\mathscr F))
\]
et comme par hypothèse $\mathscr F$ est un $\mathcal O_Y$-Module de
présentation finie,
$f^*(\mathscr H\!om_{\mathcal O_Y}(\mathscr F,\mathscr F))$ s'identifie
canoniquement à
$\mathscr H\!om_{\mathcal O_X}(f^*(\mathscr F),f^*(\mathscr F))$
$(\mathbf{0}_{\mathrm I},6.7.6)$, d'où la conclusion.

\begin{proposition}[2.1.12]
\label{IV.2.1.12-fr}
Soient $X$ un préschéma, $\mathscr F$ un $\mathcal O_X$-Module de présentation
finie, $x$ un point de $X$. Les conditions suivantes sont équivalentes~:
\begin{enumerate}
\item[\emph{a)}] $\mathscr F_x$ est un $\mathcal O_x$-module plat.
\item[\emph{b)}] Il existe un voisinage ouvert $U$ de $x$ tel que
$\mathscr F|U$ soit un $(\mathcal O_X|U)$-Module localement libre.
\end{enumerate}
\end{proposition}

En effet, $\mathscr F_x$ est un $\mathcal O_x$-module de présentation finie
et $\mathcal O_x$ un anneau local~; il revient donc au même de dire que
$\mathscr F_x$ est un $\mathcal O_x$-module plat ou un $\mathcal O_x$-module
libre (Bourbaki, \emph{Alg. comm.}, chap.~II, \S~3,
n\textsuperscript{o}~2, cor.~2 de la prop.~5)~; d'où la conclusion, compte
tenu de $(\mathbf{0}_{\mathrm I},5.2.7)$. On notera que la proposition est valable pour
un \emph{espace annelé en anneaux locaux} quelconque.

\begin{proposition}[2.1.13]
\label{IV.2.1.13-fr}
Soit $f:X\to Y$ un morphisme de préschémas. Si $f$ est plat en un point
$x\in X$ et si l'anneau $\mathcal O_x$ est réduit (\textup{resp.} intègre et
intégralement clos), alors l'anneau $\mathcal O_{f(x)}$
